\section{Geometric setting and the theorem}\label{ads:sec:setting}

Fix $m_0>0$, and let $a>0$ be the unique positive root of
$1+a^2-2m_0/a=0$. On
\[
 M=[0,\infty)\times S^2,\qquad S=\{0\}\times S^2,
\]
let
\begin{equation}\label{ads:eq:background}
 g_0=\dd s^2+r(s)^2\sigma,\qquad
 r(0)=a,\qquad r'(s)>0\quad(s>0),\qquad
 (r')^2=1+r^2-\frac{2m_0}{r},
\end{equation}
where $\sigma$ is the unit round metric. The coordinate $s$ is smooth
at the boundary, and $n=\partial_s$ points from $S$ toward infinity.
On the end, $r=r(s)$ is a coordinate and
\begin{equation}\label{ads:eq:hyperbolic-reference}
 b=\frac{\dd r^2}{1+r^2}+r^2\sigma
   =\dd\eta^2+\sinh^2\eta\,\sigma,\qquad
 \eta=\operatorname{arsinh}r.
\end{equation}
In particular, $r$ is asymptotic to a positive constant times $e^s$.

Our Laplacian is $\Delta=\Div\grad$, with nonpositive eigenvalues
on a closed manifold. The spaces $C^{k,\alpha}_\tau$ use the
hyperbolic metric $b$ in the end: the tensor and its covariant
derivatives through order $k$ are $O(e^{-\tau s})$, and the
correspondingly weighted $\alpha$-H\"older seminorms on unit
background balls are bounded. On a compact set we use ordinary
$C^{k,\alpha}$ regularity. This definition is equivalent to one
using hyperbolic distance.

Fix $0<\alpha<1$, $3/2<\tau<3$, and a smooth symmetric tensor
$q\in C^{1,\alpha}_\tau$, smooth up to $S$, satisfying
\begin{equation}\label{ads:eq:tt-seed}
 \tr_{g_0}q=0,\qquad \Div_{g_0}q=0.
\end{equation}
Throughout this part the seed $q$ is independent of $\eps$.
Suppose a smooth branch of positive solutions, defined for
$0\leq\eps<\eps_*$, satisfies
\begin{align}
 \Delta_{g_0}\phi_\eps
 &=\frac34(\phi_\eps^5-\phi_\eps)
   -\frac{\eps^2}{8}|q|_{g_0}^2\phi_\eps^{-7}
       &&\text{on }M,\label{ads:eq:conformal-equation}\\
 \partial_s\phi_\eps
 &=\frac{\eps}{4}q_{ss}\phi_\eps^{-3}
       &&\text{on }S,\label{ads:eq:conformal-boundary}\\
 \phi_0&=1,\qquad
 \phi_\eps-1\in C^{2,\alpha}_\tau,\qquad
 \|\phi_\eps-1\|_{C^{2,\alpha}_\tau}\longrightarrow0
       &&\text{as }\eps\longrightarrow0.\label{ads:eq:branch}
\end{align}
No differentiability in stronger weighted spaces is assumed.
Set
\begin{equation}\label{ads:eq:actual-data}
 g_\eps=\phi_\eps^4g_0,\qquad
 K_\eps=\eps\phi_\eps^{-2}q,\qquad
 A_\eps=|S|_{g_\eps}.
\end{equation}

\subsection{Mass and boundary conventions}

Let $x_i$, $1\leq i\leq3$, be the coordinate functions on the unit
sphere. Define $V_0=\sqrt{1+r^2}$, $V_i=rx_i$, and
$e_\eps=g_\eps-b$. We use the flux normalization
\begin{equation}\label{ads:eq:mass-flux}
\begin{split}
 p_\mu(\eps)=\frac1{16\pi}\lim_{R\to\infty}
 \int_{\{r=R\}}\big[
 &V_\mu(\Div_b e_\eps-\dd\tr_b e_\eps)\\
 &+(\tr_b e_\eps)\dd V_\mu
      -e_\eps(\grad_bV_\mu,\cdot)
 \big](\nu_b)\,\dd A_b ,
\end{split}
\end{equation}
where $\nu_b$ points toward increasing $r$. All operators and the
area form in this formula are those of $b$. Assume that these
limits are finite and that the covector is future timelike for
all sufficiently small positive $\eps$. Thus
\[
 m_{\AH}(\eps)=\sqrt{p_0(\eps)^2-\sum_{i=1}^3p_i(\eps)^2}
\]
is the positive mass norm. This normalization gives $p_0(0)=m_0$
and $p_i(0)=0$. The $p_i$ are spatial components of the hyperbolic
mass covector, not angular momentum charges.

For a two-sided surface with unit normal $\nu$ toward the chosen
end, put
\[
 H=\Div_\Sigma\nu,\qquad
 \theta_+=H+\tr_\Sigma K.
\]
For a spacetime realization our convention is
$K(X,Y)=\bar g(\bar\nabla_Xn_{\mathrm{future}},Y)$;
in Gaussian normal coordinates this means $\partial_tg=2K$.
A future marginally outer trapped surface, or MOTS, has
$\theta_+=0$.

The standard conformal identities give
\begin{equation}\label{ads:eq:constraints}
 \tr_{g_\eps}K_\eps=0,\qquad
 \Div_{g_\eps}K_\eps=0,\qquad
 R_{g_\eps}+6=|K_\eps|_{g_\eps}^2.
\end{equation}
Indeed, for a tracefree $q$ the divergence identity is
$\Div_{\phi^4g_0}(\phi^{-2}q)=\phi^{-6}\Div_{g_0}q$.
Also
\[
 R_{\phi^4g_0}
 =\phi^{-5}(-8\Delta_{g_0}\phi-6\phi)
 =-6+\eps^2\phi^{-12}|q|_{g_0}^2.
\]
The boundary $S$ is totally geodesic for $g_0$.
Its normal in $g_\eps$ is $\phi_\eps^{-2}\partial_s$, and
\begin{equation}\label{ads:eq:boundary-mots}
 H_{g_\eps}(S)=4\phi_\eps^{-3}\partial_s\phi_\eps
       =\eps\phi_\eps^{-6}q_{ss},\qquad
 \tr_SK_\eps=-\eps\phi_\eps^{-6}q_{ss}.
\end{equation}
Consequently it is a future MOTS. Its mean curvature need not vanish.

\subsection{Statement}

Define
\begin{equation}\label{ads:eq:deficit}
 b_*(A)=\sqrt{\frac A{16\pi}}\left(1+\frac A{4\pi}\right),
 \qquad
 \mathfrak D_q(\eps)=m_{\AH}(\eps)-b_*(A_\eps).
\end{equation}
A tensor is called \emph{radial} if it is invariant under every
rotation of the $S^2$ factor.

\begin{theorem}\label{ads:thm:main}
For every background, fixed TT seed, and solution branch satisfying
\eqref{ads:eq:background}--\eqref{ads:eq:branch} and the mass assumptions
above, there is $0<\eps_0\leq\eps_*$ such that
\[
 m_{\AH}(\eps)\geq
 \sqrt{\frac{A_\eps}{16\pi}}\left(1+\frac{A_\eps}{4\pi}\right)
 \qquad(0\leq\eps<\eps_0).
\]
For all sufficiently small positive $\eps$, equality holds if the
seed is radial, and the inequality is strict otherwise.
\end{theorem}

The theorem makes no sign assumption on $q_{ss}|_S$.
It also does not require outermostness or outer area-minimization.
In particular it implies the following formulation with the usual
geometric boundary hypotheses. Here an enclosing surface separates $S$
from the end and bounds with $S$ a compact region. Write $A_{\min}(S)$
for the infimum of areas of smooth enclosing surfaces, allowing
disconnected competitors and $S$ itself.

\begin{corollary}\label{ads:cor:penrose}
In addition to the hypotheses of Theorem~\ref{ads:thm:main}, suppose
$q_{ss}\geq0$ on $S$. Assume, for every sufficiently small positive
$\eps$, that no compact smooth embedded two-sided surface in the
interior enclosing $S$ has $\theta_+\leq0$ everywhere, and that
every smooth enclosing surface has area at least $A_\eps$.
Disconnected enclosing competitors are allowed, and $S$ itself is
allowed in the area infimum. Then the exact local Penrose inequality
holds with $A_{\min}(S)=A_\eps$.
\end{corollary}

The outermostness condition excludes
all weakly future outer trapped enclosing surfaces, not only
additional MOTS. These additional assumptions are relevant to the
geometric formulation of the conjecture but are not needed in the
coefficient argument below.

\subsection{Static operators and Taylor notation}\label{ads:sec:taylor}

We use
\begin{equation}\label{ads:eq:static-operators}
 \begin{gathered}
 N=r',\qquad
 \kappa=N'(0)=\frac{1+3a^2}{2a},\qquad
 B=\Ric_{g_0}+3g_0,\\
 L=\Delta_{g_0}-3,\qquad
 D=-\Delta_{a^2\sigma}+\frac{2\kappa}{a}.
 \end{gathered}
\end{equation}
The letter $D$ denotes the boundary operator, while
$\mathfrak D_q$ denotes the mass deficit.

\begin{lemma}\label{ads:lem:static}
The tensor $B$ and the static lapse satisfy
\begin{gather}
 B=B_s\,\dd s^2+B_t r^2\sigma,\qquad
 B_s=1-\frac{2m_0}{r^3},\qquad
 B_t=1+\frac{m_0}{r^3},\label{ads:eq:B-components}\\
 \nabla^2N=NB,\qquad
 \tr_{g_0}B=3,\qquad \Div_{g_0}B=0,\qquad
 LN=0.\label{ads:eq:static-identities}
\end{gather}
Moreover $D$ is positive on all spherical harmonics and
\begin{equation}\label{ads:eq:boundary-constants}
 b_*(4\pi a^2)=m_0,\qquad
 b_*'(4\pi a^2)=\frac{\kappa}{8\pi}.
\end{equation}
\end{lemma}

\begin{proof}
The radial and tangential sectional curvatures of
$\dd s^2+r^2\sigma$ are $-r''/r$ and $(1-(r')^2)/r^2$.
Differentiating \eqref{ads:eq:background} gives
$N'=r+m_0/r^2$ and $N''=N(1-2m_0/r^3)$.
The Ricci eigenvalues are therefore $-2-2m_0/r^3$ and
$-2+m_0/r^3$, giving \eqref{ads:eq:B-components} and $R_{g_0}=-6$.
The radial and tangential Hessian eigenvalues of $N$ are
$N''$ and $N'N/r$, respectively, proving the first identity
in \eqref{ads:eq:static-identities}. Its trace gives $LN=0$.
The divergence identity follows from the contracted Bianchi
identity and the constant scalar curvature.
At $S$, $B_t=\kappa/a$ and $2\kappa/a>0$, so $D>0$.
Finally $2m_0=a(1+a^2)$ gives \eqref{ads:eq:boundary-constants}
by direct differentiation.
\end{proof}

We write $[\mathcal F]_j$ for the coefficient of $\eps^j$ in a
Taylor expansion; coefficients do not include an extra factorial.
Thus $u_j[q]=[\phi_\eps-1]_j$ and $c_j[q]=[\mathfrak D_q]_j$
when the indicated actual expansions have been established.
For a tensor $T$ we also use the weighted function norm
\[
 \|T\|_\beta=\sup_M e^{\beta s}|T|_{g_0}.
\]
On the end this is equivalent to using $b$.
Finite angular type means that the span of all rotational
pullbacks of the function or tensor is finite-dimensional.
For functions this is precisely a finite sum of spherical-harmonic
spaces, with arbitrary smooth radial coefficients.

Some geometric arguments use Taylor families of metrics and
surfaces only to order $p$, with identities read in the finite Taylor
algebra $\mathbb R[\eps]/(\eps^{p+1})$. On any compact radial interval,
their coefficients can be represented by a smooth family and
all differential-geometric identities then hold to the retained
order. At $S$, a formal displacement may have negative first
coefficient. It is evaluated using the given smooth one-sided
jets: extend sufficiently many coefficients smoothly across
$s=0$, perform the finite Taylor calculation, and retain only
the prescribed order. Every boundary derivative of an identity
holding on $s\geq0$ is fixed by those jets. This convention
does not assume a constrained extension at a fixed point with
$s<0$. The time-recursion argument below will specify how the
constraint identities are used after such a displacement.
